\frontchapter{Abstract}
Approximate nearest-neighbour (ANN) search supports similarity retrieval
over large collections of embedding vectors. GPUs accelerate this search,
but limited device memory makes compressed indexes attractive. JHQ is a
hierarchical quantisation method that screens candidates with compact
codes and refines a smaller set using residual codes. This thesis develops
JHQ-GPU to study how this representation can support efficient GPU
construction and search.

The implementation combines GPU model fitting and batched encoding with
bounded input staging and a packed, list-ordered layout. Its scan replaces
each 256-entry distance table with two 16-entry tables, preserving scores
in real arithmetic. Grouped scanning preserves exact global candidate
selection under a fixed ordering of computed scores and positions. Calibration chooses how
many candidates to refine without ground-truth labels, by comparing
returned neighbours with a larger-budget reference.

Evaluation on an NVIDIA RTX~5090 compares methods on Vogue, arXiv,
OpenAI-1536 and OpenAI-3072. At recall 0.97 and batch 1,000, JHQ-GPU
delivers $1.32$--$1.75\times$ the host-to-host throughput of the CAGRA fp32 graph index
across all four datasets and is the fastest tested method on three.
For the larger BGE and Stella text-embedding datasets, whose raw FP32 vectors exceed
device memory, JHQ-GPU enables single-GPU search; the largest index
contains 17.78 million vectors and uses 20.44\,GiB of resident GPU memory.
Using 500 unlabelled calibration queries keeps test-query recall within
$10^{-3}$ of the best tested budget in all twelve dataset--probe settings.
